% GENERATED FILE. Edit canonical lessons, then run npm run generate:books.
% canonical-source-hash: fnv1a64:380e6953036ca13b
% canonical-lessons: GE-C145-benutzen, GE-C145-einkaufen, GE-C145-einladen, GE-C145-empfehlen, GE-C145-entscheiden

\chapter{To use, To go shopping, To invite, To recommend, To decide}
\label{ch:ge-to-use-to-go-shopping-to-invite-to-recommend-to-decide}
\hlchaptermodality{145}

\begin{chapteropening}
\textbf{By the end of this chapter:} \emph{I can say to use, to go shopping, to invite, to recommend and to decide.}

Complete the last lesson of chapter 145: I can say to use, to go shopping, to invite, to recommend and to decide.
\end{chapteropening}

\section[benutzen]{benutzen — to use}

\label{lesson:GE-C145-benutzen}



\begin{quote}
Before the new one: say the German for to go out, then the German for to look.
\end{quote}

\subsection*{You'll want to know: benutzen}
\textbf{benutzen} — \textquotedblleft{}to use\textquotedblright{}.

\begin{cousinweb}
\textbf{benutzen} is \emph{be-} + the verb \textquotedblleft{}to use, to profit from\textquotedblright{}, Old High German \emph{nuzzōn} $\leftarrow$ Proto-Germanic \emph{*nutōną} (EtymDB 2.1).
\end{cousinweb}

\begin{grammarlens}[title={What you've built}]
One more word for this chapter.
\end{grammarlens}

\subsection*{Your turn}
\begin{itemize}
\raggedright
  \item \emph{Say it:} \emph{benutzen}
  \item \emph{Say it:} \emph{benutzen}, once more
  \item \emph{Say it:} say \emph{benutzen}
  \item \emph{Recall:} say the German for to go out, then the German for to look, then say \emph{benutzen} again
\end{itemize}

\begin{tcolorbox}[breakable,colback=teal!4,colframe=teal!35!black,title={Before you move on}]
What does benutzen mean? (\textquotedblleft{}To use\textquotedblright{}.) Say it once more.
\end{tcolorbox}
\section[einkaufen]{einkaufen — to go shopping}

\label{lesson:GE-C145-einkaufen}



\begin{quote}
Before the new one: say the German for to look, then the German for to use.
\end{quote}

\subsection*{You'll want to know: einkaufen}
\textbf{einkaufen} — \textquotedblleft{}to go shopping\textquotedblright{}.

\begin{cousinweb}
\textbf{einkaufen} is \emph{ein-}, \textquotedblleft{}in\textquotedblright{}, + \emph{kaufen}, \textquotedblleft{}to buy\textquotedblright{}: to buy in supplies. \emph{kaufen} goes back to Latin \emph{caupō}, \textquotedblleft{}tradesman\textquotedblright{}. (EtymDB 2.1)
\end{cousinweb}

\begin{grammarlens}[title={What you've built}]
One more word for this chapter.
\end{grammarlens}

\subsection*{Your turn}
\begin{itemize}
\raggedright
  \item \emph{Say it:} \emph{einkaufen}
  \item \emph{Say it:} \emph{einkaufen}, once more
  \item \emph{Say it:} say \emph{einkaufen}
  \item \emph{Recall:} say the German for to look, then the German for to use, then say \emph{einkaufen} again
\end{itemize}

\begin{tcolorbox}[breakable,colback=teal!4,colframe=teal!35!black,title={Before you move on}]
What does einkaufen mean? (\textquotedblleft{}To go shopping\textquotedblright{}.) Say it once more.
\end{tcolorbox}
\section[einladen]{einladen — to invite}

\label{lesson:GE-C145-einladen}



\begin{quote}
Before the new one: say the German for to use, then the German for to go shopping.
\end{quote}

\subsection*{You'll want to know: einladen}
\textbf{einladen} — \textquotedblleft{}to invite\textquotedblright{}.

\begin{cousinweb}
\textbf{einladen} $\leftarrow$ Old High German \emph{inladōn} (EtymDB 2.1): \emph{ein-}, \textquotedblleft{}in\textquotedblright{}, + an old verb meaning \textquotedblleft{}to call, to summon\textquotedblright{}. It is a different word from \emph{Laden}, \textquotedblleft{}shop\textquotedblright{}.
\end{cousinweb}

\begin{grammarlens}[title={What you've built}]
One more word for this chapter.
\end{grammarlens}

\subsection*{Your turn}
\begin{itemize}
\raggedright
  \item \emph{Say it:} \emph{einladen}
  \item \emph{Say it:} \emph{einladen}, once more
  \item \emph{Say it:} say \emph{einladen}
  \item \emph{Recall:} say the German for to use, then the German for to go shopping, then say \emph{einladen} again
\end{itemize}

\begin{tcolorbox}[breakable,colback=teal!4,colframe=teal!35!black,title={Before you move on}]
What does einladen mean? (\textquotedblleft{}To invite\textquotedblright{}.) Say it once more.
\end{tcolorbox}
\section[empfehlen]{empfehlen — to recommend}

\label{lesson:GE-C145-empfehlen}



\begin{quote}
Before the new one: say the German for to go shopping, then the German for to invite.
\end{quote}

\subsection*{You'll want to know: empfehlen}
\textbf{empfehlen} — \textquotedblleft{}to recommend\textquotedblright{}.

\begin{cousinweb}
\textbf{empfehlen} $\leftarrow$ Middle High German \emph{enpfelhen}, \textquotedblleft{}to entrust to someone\textquotedblright{} (EtymDB 2.1) — inherited, built from \emph{ent-} and an old verb for \textquotedblleft{}to hand over\textquotedblright{}. To recommend is to entrust.
\end{cousinweb}

\begin{grammarlens}[title={What you've built}]
One more word for this chapter.
\end{grammarlens}

\subsection*{Your turn}
\begin{itemize}
\raggedright
  \item \emph{Say it:} \emph{empfehlen}
  \item \emph{Say it:} \emph{empfehlen}, once more
  \item \emph{Say it:} say \emph{empfehlen}
  \item \emph{Recall:} say the German for to go shopping, then the German for to invite, then say \emph{empfehlen} again
\end{itemize}

\begin{tcolorbox}[breakable,colback=teal!4,colframe=teal!35!black,title={Before you move on}]
What does empfehlen mean? (\textquotedblleft{}To recommend\textquotedblright{}.) Say it once more.
\end{tcolorbox}
\section[entscheiden]{entscheiden — to decide}

\label{lesson:GE-C145-entscheiden}



\begin{quote}
Before the new one: say the German for to invite, then the German for to recommend.
\end{quote}

\subsection*{You'll want to know: entscheiden}
\textbf{entscheiden} — \textquotedblleft{}to decide\textquotedblright{}.

\begin{cousinweb}
\textbf{entscheiden} is \emph{ent-} + the verb \textquotedblleft{}to separate, to divide\textquotedblright{}, Old High German \emph{sceidan} $\leftarrow$ Proto-Germanic \emph{*skaiþaną} (EtymDB 2.1): to decide is to cut between choices.
\end{cousinweb}

\begin{grammarlens}[title={What you've built}]
One more word for this chapter.
\end{grammarlens}

\subsection*{Your turn}
\begin{itemize}
\raggedright
  \item \emph{Say it:} \emph{entscheiden}
  \item \emph{Say it:} \emph{entscheiden}, once more
  \item \emph{Say it:} say \emph{entscheiden}
  \item \emph{Recall:} say the German for to invite, then the German for to recommend, then say \emph{entscheiden} again
\end{itemize}

\begin{tcolorbox}[breakable,colback=teal!4,colframe=teal!35!black,title={Before you move on}]
What does entscheiden mean? (\textquotedblleft{}To decide\textquotedblright{}.) Say it once more.
\end{tcolorbox}
